
\begin{answerBox}{2in}          % please don't edit this line, doing so will throw off the page layout

  %%%%%%%%%%%%%%%%%%%%%%%%%%%%%%%%%%%%%%%%%%%%%%%%%%%%%%
  %%%%%%%%%%%%%%%%%%%%%%%%%%%%%%%%%%%%%%%%%%%%%%%%%%%%%%
  %%
  %% PROBLEM 2 COLLABORATORS AND CITATIONS 
  %%
  %%%%%%%%%%%%%%%%%%%%%%%%%%%%%%%%%%%%%%%%%%%%%%%%%%%%%%
  %%%%%%%%%%%%%%%%%%%%%%%%%%%%%%%%%%%%%%%%%%%%%%%%%%%%%%
  
  \paragraph{Problem \arabic{problem} collaboration and citations}~\GradescopeComment
  {This should be on page 4 (bottom) of your PDF.}

  I collaborated on this problem with the following people:
  \begin{blue}
    \begin{itemize}
      
    \item \REPLACEME            % write "none"if none
      
    \end{itemize}
  \end{blue}

  My answers use ideas from the following sources (other than the lecture notes and textbooks): 
  \begin{blue}
    \begin{itemize}
      
    \item \REPLACEME            % write "none"if none
      
    \end{itemize}
  \end{blue}

\end{answerBox}

\newpage

\paragraph{Problem 2 answer}\GradescopeComment{This should be on page 5 (top) of your PDF.}

\begin{enumerate}[label=(\alph*)]

\item 
  %%%%%%%%%%%%%%%%%%%%%%%%%%%%%%%%%%%%%%%%%%%%%%%%%%%%%% 
  %%%%%%%%%%%%%%%%%%%%%%%%%%%%%%%%%%%%%%%%%%%%%%%%%%%%%% 
  %% 
  %% PROBLEM 2(a) ANSWER 
  %% 
  %%%%%%%%%%%%%%%%%%%%%%%%%%%%%%%%%%%%%%%%%%%%%%%%%%%%%% 
  %%%%%%%%%%%%%%%%%%%%%%%%%%%%%%%%%%%%%%%%%%%%%%%%%%%%%% 

  % Per Section 7.4,
  % define the subproblems that your algorithm solves for a given input $(v=(v_1,\ldots, v_n), T)$.

  Here are the subproblems that the algorithm solves.
  
  \begin{blue}
    For each
    \REPLACEME,
    define
    \REPLACEME
    to be
    \REPLACEME.                 % explain in words
    That is, 
    \[\REPLACEME = \REPLACEME.\] % clarify it with a mathematical definition
  \end{blue}

\item
  %%%%%%%%%%%%%%%%%%%%%%%%%%%%%%%%%%%%%%%%%%%%%%%%%%%%%% 
  %%%%%%%%%%%%%%%%%%%%%%%%%%%%%%%%%%%%%%%%%%%%%%%%%%%%%% 
  %% 
  %% PROBLEM 2(b) ANSWER 
  %% 
  %%%%%%%%%%%%%%%%%%%%%%%%%%%%%%%%%%%%%%%%%%%%%%%%%%%%%% 
  %%%%%%%%%%%%%%%%%%%%%%%%%%%%%%%%%%%%%%%%%%%%%%%%%%%%%% 

 % Which subproblem gives the final answer?

  The correct output for the given input is
  \[
    \begin{blue}
      \REPLACEME.
    \end{blue}
  \]

\item
  %%%%%%%%%%%%%%%%%%%%%%%%%%%%%%%%%%%%%%%%%%%%%%%%%%%%%% 
  %%%%%%%%%%%%%%%%%%%%%%%%%%%%%%%%%%%%%%%%%%%%%%%%%%%%%% 
  %% 
  %% PROBLEM 2(c) ANSWER 
  %% 
  %%%%%%%%%%%%%%%%%%%%%%%%%%%%%%%%%%%%%%%%%%%%%%%%%%%%%% 
  %%%%%%%%%%%%%%%%%%%%%%%%%%%%%%%%%%%%%%%%%%%%%%%%%%%%%% 

  % State an appropriate recurrence relation, including boundary cases. 

  Here is the recurrence relation:

  \begin{blue}
    \[
      \REPLACEME =
      \begin{cases}
        \REPLACEME & \text{ if } \REPLACEME, \\
        \REPLACEME & \text{ if } \REPLACEME, \\
        \REPLACEME & \text{ if } \REPLACEME, \\
        \REPLACEME & \text{ if } \REPLACEME. \\
      \end{cases}
    \]
  \end{blue}

\item
  %%%%%%%%%%%%%%%%%%%%%%%%%%%%%%%%%%%%%%%%%%%%%%%%%%%%%% 
  %%%%%%%%%%%%%%%%%%%%%%%%%%%%%%%%%%%%%%%%%%%%%%%%%%%%%% 
  %% 
  %% PROBLEM 2(d) ANSWER 
  %% 
  %%%%%%%%%%%%%%%%%%%%%%%%%%%%%%%%%%%%%%%%%%%%%%%%%%%%%% 
  %%%%%%%%%%%%%%%%%%%%%%%%%%%%%%%%%%%%%%%%%%%%%%%%%%%%%% 

  % In terms of $n$ and $T$, what is the big-$\Theta$ running time of the resulting iterative (or recursive and memoized) algorithm?

  The running time is

  \begin{blue}
    \[
      \REPLACEME.
    \]
  \end{blue}
  
\end{enumerate}

%% THE COLLABORATORS AND CITATIONS STATEMENT IS AT THE TOP OF THIS FILE

%%% Local Variables:
%%% mode: latex
%%% TeX-master: "homework_9"
%%% End:
